% Fichier engendré par tools/generate-tex.py à partir de 06-statistiques-et-probabilites.
% Ne pas éditer à la main : tout le texte provient des commentaires du fichier Lean.
\documentclass[11pt,a4paper]{article}

% Compile aussi bien avec pdflatex qu'avec xelatex ou tectonic.
\usepackage{iftex}
\ifPDFTeX
  \usepackage[T1]{fontenc}
  \usepackage[utf8]{inputenc}
\else
  \usepackage{fontspec}
\fi
\usepackage[french]{babel}
\usepackage{amsmath,amssymb,amsthm}
\usepackage[margin=2.5cm]{geometry}
\usepackage[hidelinks]{hyperref}

\theoremstyle{definition}
\newtheorem{definition}{Définition}
\newtheorem{exemple}{Exemple}
\theoremstyle{plain}
\newtheorem{theoreme}{Théorème}
\newtheorem{lemme}[theoreme]{Lemme}

% Renvoi à la déclaration Lean, une ligne après la démonstration, aligné à droite.
\newcommand{\source}[2]{\par\smallskip\noindent\hfill{\footnotesize\href{#1}{\texttt{#2}}}\par\smallskip}

\title{Statistiques et probabilités}
\date{}

\begin{document}
\maketitle
% Les identifiants Lean en machine à écrire ne se coupent pas : on tolère des
% espacements plus lâches plutôt que des lignes qui débordent.
\sloppy

\section{StatistiquesEtProbabilites}

\noindent
Collège \textemdash{} section \og{} Statistiques et probabilités \fg{}.
Une série statistique est une liste de nombres, une expérience aléatoire est donnée par
la probabilité de chacune de ses issues.
Énoncés et démonstrations en français : voir StatistiquesEtProbabilites.tex.

\subsection{Définitions}

\noindent
Toutes les définitions du fichier sont posées ici, avant les démonstrations.

\begin{definition}
La \emph{moyenne} d'une série $s = (x_1, \dots, x_n)$ est
\[ \bar{x} = \frac{x_1 + \cdots + x_n}{n} . \]
\end{definition}
\source{https://github.com/Commutator-IO/learn-lean/blob/main/courses/01-college/06-statistiques-et-probabilites/StatistiquesEtProbabilites.lean\#L20}{StatistiquesEtProbabilites.lean\#L20}

\begin{definition}
La série \emph{translatée} de $c$ est celle obtenue en ajoutant $c$ à chaque valeur.
\end{definition}
\source{https://github.com/Commutator-IO/learn-lean/blob/main/courses/01-college/06-statistiques-et-probabilites/StatistiquesEtProbabilites.lean\#L23}{StatistiquesEtProbabilites.lean\#L23}

\begin{definition}
La \emph{moyenne pondérée} de valeurs $v_i$ affectées d'effectifs $n_i$ est
\[ \frac{\sum_i n_i v_i}{\sum_i n_i} . \]
\end{definition}
\source{https://github.com/Commutator-IO/learn-lean/blob/main/courses/01-college/06-statistiques-et-probabilites/StatistiquesEtProbabilites.lean\#L26}{StatistiquesEtProbabilites.lean\#L26}

\begin{definition}
$m$ est une \emph{médiane} de la série $s$ lorsque au moins la moitié des valeurs lui
sont inférieures ou égales, et au moins la moitié lui sont supérieures ou égales.
\end{definition}
\source{https://github.com/Commutator-IO/learn-lean/blob/main/courses/01-college/06-statistiques-et-probabilites/StatistiquesEtProbabilites.lean\#L31}{StatistiquesEtProbabilites.lean\#L31}

\begin{definition}
Une \emph{loi de probabilité} sur un ensemble fini d'issues $\Omega$ est la donnée
d'un nombre $p(\omega) \ge 0$ pour chaque issue, tel que
\[ \sum_{\omega \in \Omega} p(\omega) = 1 . \]
\end{definition}
\source{https://github.com/Commutator-IO/learn-lean/blob/main/courses/01-college/06-statistiques-et-probabilites/StatistiquesEtProbabilites.lean\#L36}{StatistiquesEtProbabilites.lean\#L36}

\begin{definition}
La \emph{probabilité} d'un événement $A$ est la somme des probabilités de ses issues :
$P(A) = \sum_{\omega \in A} p(\omega)$.
\end{definition}
\source{https://github.com/Commutator-IO/learn-lean/blob/main/courses/01-college/06-statistiques-et-probabilites/StatistiquesEtProbabilites.lean\#L42}{StatistiquesEtProbabilites.lean\#L42}

\subsection{Moyenne : linéarité et moyenne pondérée}

\begin{theoreme}
Ajouter un même nombre à toutes les valeurs ajoute ce nombre à la moyenne :
\[ \overline{x + c} = \bar{x} + c . \]
\end{theoreme}

\begin{proof}
La somme des valeurs translatées vaut la somme des valeurs augmentée de $nc$ :
\[ \frac{\sum_i (x_i + c)}{n} = \frac{\left(\sum_i x_i\right) + nc}{n}
   = \bar{x} + c , \]
la division par $n$ étant licite puisque la série n'est pas vide.
\end{proof}
\source{https://github.com/Commutator-IO/learn-lean/blob/main/courses/01-college/06-statistiques-et-probabilites/StatistiquesEtProbabilites.lean\#L48}{StatistiquesEtProbabilites.lean\#L48}

\begin{theoreme}
Quand tous les effectifs valent un, la moyenne pondérée est la moyenne ordinaire :
\[ \frac{\sum_i 1 \times v_i}{\sum_i 1} = \frac{\sum_i v_i}{k} . \]
\end{theoreme}

\begin{proof}
Le dénominateur $\sum_i 1$ vaut le nombre $k$ de valeurs, et chaque terme du numérateur
se simplifie. La moyenne pondérée n'est donc que la moyenne, écrite sans répéter les
valeurs identiques.
\end{proof}
\source{https://github.com/Commutator-IO/learn-lean/blob/main/courses/01-college/06-statistiques-et-probabilites/StatistiquesEtProbabilites.lean\#L60}{StatistiquesEtProbabilites.lean\#L60}

\subsection{La moyenne est comprise entre le minimum et le maximum}

\begin{theoreme}
La moyenne est comprise entre deux bornes qui encadrent toutes les valeurs : si
$a \le x_i \le b$ pour tout $i$, alors $a \le \bar{x} \le b$.
\end{theoreme}

\begin{proof}
Chaque valeur étant au moins $a$, la somme est au moins $na$, donc $\bar{x} \ge a$
après division par $n > 0$ :
\[ na = \sum_i a \le \sum_i x_i . \]
Le même raisonnement avec $b$ en majorant donne l'autre inégalité. En particulier, la
moyenne est comprise entre le minimum et le maximum de la série.
\end{proof}
\source{https://github.com/Commutator-IO/learn-lean/blob/main/courses/01-college/06-statistiques-et-probabilites/StatistiquesEtProbabilites.lean\#L69}{StatistiquesEtProbabilites.lean\#L69}

\subsection{La moyenne de plusieurs moyennes n'est pas la moyenne de la série globale}

\begin{exemple}
La moyenne de plusieurs moyennes n'est pas la moyenne de la série globale :
\[ \overline{\bigl(\overline{(0,0,0)},\ \overline{(6)}\bigr)} \ne \overline{(0,0,0,6)} . \]
\end{exemple}

\begin{proof}
Les deux groupes ont pour moyennes $0$ et $6$, dont la moyenne vaut $3$. La série
réunie, elle, a pour moyenne $\tfrac{0+0+0+6}{4} = 1{,}5$. Moyenner des moyennes revient
à donner le même poids à des groupes de tailles différentes.
\end{proof}
\source{https://github.com/Commutator-IO/learn-lean/blob/main/courses/01-college/06-statistiques-et-probabilites/StatistiquesEtProbabilites.lean\#L90}{StatistiquesEtProbabilites.lean\#L90}

\subsection{Médiane : au moins la moitié des valeurs lui sont inférieures ou égales}

\begin{exemple}
Sur la série $(1, 2, 3, 4, 5)$, la valeur centrale $3$ est bien une médiane.
\end{exemple}

\begin{proof}
Trois valeurs sur cinq lui sont inférieures ou égales ($1$, $2$, $3$) et trois lui sont
supérieures ou égales ($3$, $4$, $5$) : dans les deux cas, plus de la moitié des cinq
valeurs.
\end{proof}
\source{https://github.com/Commutator-IO/learn-lean/blob/main/courses/01-college/06-statistiques-et-probabilites/StatistiquesEtProbabilites.lean\#L98}{StatistiquesEtProbabilites.lean\#L98}

\begin{theoreme}
Une médiane sépare la série en deux moitiés : au moins la moitié des valeurs lui sont
inférieures ou égales.
\end{theoreme}

\begin{proof}
C'est la première moitié de la définition, isolée pour l'usage.
\end{proof}
\source{https://github.com/Commutator-IO/learn-lean/blob/main/courses/01-college/06-statistiques-et-probabilites/StatistiquesEtProbabilites.lean\#L103}{StatistiquesEtProbabilites.lean\#L103}

\subsection{Probabilité : \texttt{0 \(\le\) P(A) \(\le\) 1} et somme des probabilités des issues = 1}

\begin{theoreme}
Une probabilité est positive : $P(A) \ge 0$.
\end{theoreme}

\begin{proof}
C'est une somme de nombres positifs, chacun étant la probabilité d'une issue.
\end{proof}
\source{https://github.com/Commutator-IO/learn-lean/blob/main/courses/01-college/06-statistiques-et-probabilites/StatistiquesEtProbabilites.lean\#L111}{StatistiquesEtProbabilites.lean\#L111}

\begin{theoreme}
Une probabilité ne dépasse pas un : $P(A) \le 1$.
\end{theoreme}

\begin{proof}
La somme sur $A$ est majorée par la somme sur toutes les issues, puisqu'on ajoute des
termes positifs ; or cette dernière vaut $1$ par définition d'une loi.
\end{proof}
\source{https://github.com/Commutator-IO/learn-lean/blob/main/courses/01-college/06-statistiques-et-probabilites/StatistiquesEtProbabilites.lean\#L115}{StatistiquesEtProbabilites.lean\#L115}

\begin{theoreme}
La probabilité de l'événement certain vaut un : $P(\Omega) = 1$.
\end{theoreme}

\begin{proof}
C'est exactement la condition imposée à une loi de probabilité : la somme des
probabilités des issues vaut un.
\end{proof}
\source{https://github.com/Commutator-IO/learn-lean/blob/main/courses/01-college/06-statistiques-et-probabilites/StatistiquesEtProbabilites.lean\#L122}{StatistiquesEtProbabilites.lean\#L122}

\subsection{Événement contraire : \texttt{P(Ā) = 1 − P(A)}}

\begin{theoreme}
Événement contraire : $P(\bar{A}) = 1 - P(A)$.
\end{theoreme}

\begin{proof}
Les issues de $A$ et celles de $\bar{A}$ forment ensemble toutes les issues, sans
recouvrement :
\[ P(A) + P(\bar{A}) = \sum_{\omega \in \Omega} p(\omega) = 1 , \]
d'où le résultat en retranchant $P(A)$.
\end{proof}
\source{https://github.com/Commutator-IO/learn-lean/blob/main/courses/01-college/06-statistiques-et-probabilites/StatistiquesEtProbabilites.lean\#L127}{StatistiquesEtProbabilites.lean\#L127}

\subsection{Équiprobabilité : \texttt{P(A) = card(A) / card(Ω)}}

\begin{theoreme}
Équiprobabilité : si toutes les issues ont la même probabilité, alors
\[ P(A) = \frac{\operatorname{card} A}{\operatorname{card} \Omega} . \]
\end{theoreme}

\begin{proof}
Chaque issue ayant la probabilité $1/\operatorname{card}\Omega$, la somme sur $A$
compte $\operatorname{card} A$ fois cette valeur :
\[ P(A) = \operatorname{card} A \times \frac{1}{\operatorname{card}\Omega} . \]
C'est la règle « nombre de cas favorables sur nombre de cas possibles ».
\end{proof}
\source{https://github.com/Commutator-IO/learn-lean/blob/main/courses/01-college/06-statistiques-et-probabilites/StatistiquesEtProbabilites.lean\#L137}{StatistiquesEtProbabilites.lean\#L137}

\subsection{Expérience à deux épreuves : produit des probabilités le long d'une branche}

\begin{definition}
La \emph{loi produit} de deux épreuves indépendantes affecte à un couple d'issues le
produit de leurs probabilités : $p(a, b) = p(a)\,p'(b)$. C'en est bien une : ses valeurs
sont positives et leur somme vaut un.
\end{definition}
\source{https://github.com/Commutator-IO/learn-lean/blob/main/courses/01-college/06-statistiques-et-probabilites/StatistiquesEtProbabilites.lean\#L146}{StatistiquesEtProbabilites.lean\#L146}

\begin{theoreme}
Le long d'une branche de l'arbre, les probabilités se multiplient :
$P\bigl((a, b)\bigr) = P(a) \times P(b)$.
\end{theoreme}

\begin{proof}
C'est la définition de la loi produit. Ce qui demande une vérification, c'est qu'elle
soit une loi : la somme de tous les produits vaut
\[ \sum_{a}\sum_{b} p(a)p'(b) = \sum_a p(a) \left(\sum_b p'(b)\right)
   = \sum_a p(a) \times 1 = 1 . \]
\end{proof}
\source{https://github.com/Commutator-IO/learn-lean/blob/main/courses/01-college/06-statistiques-et-probabilites/StatistiquesEtProbabilites.lean\#L156}{StatistiquesEtProbabilites.lean\#L156}

\subsection{Fréquence observée et probabilité : stabilisation quand \texttt{n} grandit}

\begin{theoreme}
La fluctuation d'échantillonnage s'amenuise : la suite $1/\sqrt{n}$, qui donne l'ordre
de grandeur de l'écart entre fréquence observée et probabilité, tend vers zéro.
\end{theoreme}

\begin{proof}
La racine carrée tend vers l'infini avec $n$, donc son inverse tend vers zéro.

L'énoncé du collège — « les fréquences se stabilisent autour de la probabilité quand le
nombre d'essais grandit » — est plus fort : c'est la loi des grands nombres, dont la
démonstration relève de la théorie de la mesure. Ce qui est établi ici est la décroissance
de l'ordre de grandeur $1/\sqrt{n}$ de la fluctuation, pas la convergence des fréquences
elles-mêmes.
\end{proof}
\source{https://github.com/Commutator-IO/learn-lean/blob/main/courses/01-college/06-statistiques-et-probabilites/StatistiquesEtProbabilites.lean\#L164}{StatistiquesEtProbabilites.lean\#L164}

\vfill
\noindent\rule{0.35\linewidth}{0.4pt}\par
\noindent{\footnotesize Leonardo de Moura et Sebastian Ullrich,
\emph{The Lean~4 Theorem Prover and Programming Language},
\emph{Automated Deduction — CADE~28}, Springer, 2021, p.~625--635,
\href{https://doi.org/10.1007/978-3-030-79876-5_37}{doi:10.1007/978-3-030-79876-5\_37}.\par}

\end{document}
